\begin{proposition}[Integral homology of the local degeneration]\hypertarget{tag-02X2}{}\label{night:P4.3}\label{stmt:P4.3}
\StatementRecord{night:P4.3}{proved}{02X2}{the definitions and constructions stated in the proof}

For the degeneration $U$ of Proposition~\ref{night:P4.1}, order the
regular fibre lattice as $(a,c,b,d)$, where $(a,c)$ are the cycles
$(1,i)$ in $C_i$. The regular-fibre inclusions induce the quotient maps
\[
 H_1(A)=\Lambda\longrightarrow\Lambda/\Z a=H_1(U),\qquad
 H_2(A)=\bigwedge^2\Lambda\longrightarrow
 \bigwedge^2\Lambda/\Z(a\wedge c)=H_2(U).
\]
Both target groups are free. The meridian and every product of the
meridian with a fibre one-cycle map to zero in these degrees.
\end{proposition}
\begin{proof}

\begin{enumerate}[label={$\langle1\rangle$\arabic*.},leftmargin=*]
\item A Hermitian metric gives a diffeomorphism
$L^\times/q^\Z\cong S(L)\times S^1_c$, with factors the circle bundle and
the logarithm of the norm modulo $\log q$. We first specify a homotopy
model for $V$. Blowing up a four-ball replaces it by the disc bundle of
$\OO(-1)$ along its Hopf boundary. The complement of the ball in
$B\times\Delta$ has homotopy type $B\vee S^3$; the boundary sphere is
the $S^3$ summand, since its map to $B$ is nullhomotopic. The resulting
homotopy pushout is $B\vee S^2$. Equivalently, on simply connected
covers the replacement adds one sphere at each lift of the centre and
no other homology; the comparison is an isomorphism on $\pi_1$ and on
the homology of those covers, and hence a homotopy equivalence. We may
choose the $B$ summand to be a regular $B_t$ and the sphere to be $R$,
joining their base points by a path.

\item The circle Chern class is zero on $B_t$ and $-1$ on $R$. Therefore
$S(L)$ is represented by the homotopy pushout
\[
 T^3_{(a,b,d)}\longleftarrow S^1_a\longrightarrow S^3.
\]
The second map is a Hopf fibre and is nullhomotopic. Thus this pushout
has the homotopy type $(T^3\cup_aD^2)\vee S^3$. We multiply by $S^1_c$.
The resulting integral cellular complex starts with the exterior-cell
complex of $T^4$, ordered by $(a,c,b,d)$, and adds cells $r$ in degree
two, $rc,h$ in degree three, and $hc$ in degree four. Its only nonzero
differentials are
\[
 \partial r=a,\qquad \partial(rc)=a\wedge c.
\]
We choose the orientation of $r$ to give the first equality and use the
product orientation for $rc$. The pairs $(r,a)$ and $(rc,a\wedge c)$
are integral contractible summands: the degree-one homotopy sends
$a$ to $r$ and $a\wedge c$ to $rc$, and is zero on the remaining
basis cells. Deleting these pairs leaves the zero differential and
free homology of ranks $(1,3,5,5,2)$. The inclusion of the torus
subcomplex kills exactly $a$ in degree one and $a\wedge c$ in
degree two, giving the asserted quotient maps.

\item We represent a base point, and both base generators $b,d$, away from
$0\in B$. Their products with the disc miss $R$, so the canonical section
trivializes $L$ over the resulting chains. The meridian and the products
with $b,d$ consequently bound. The products with $a,c$ bound over one
such base point. This proves all the remaining assertions on the maps.

\end{enumerate}
\end{proof}